% ===== Document Preamble =====
\documentclass[11pt]{article}
\usepackage[margin=0.5in]{geometry}
\usepackage{amsmath,amssymb}
\usepackage{xcolor}
\usepackage{enumitem}
\usepackage{needspace}
\usepackage{comment}

% ===== Colors =====
\definecolor{AnsRed}{RGB}{170,0,0}
\definecolor{WrongBlue}{RGB}{0,0,160}
\definecolor{SkipGray}{RGB}{120,120,120}

% ===== Choice list styling (tight, consistent) =====
\setlist[enumerate]{%
  leftmargin=1.5em,
  itemsep=0.25em,
  topsep=0.25em,
  label=\textbf{\Alph*.}
}

% ===== Style Wrappers =====
\newcommand{\QuestionStyle}[1]{{\large\bfseries\boldmath #1}}
\newcommand{\AnswerStyle}[1]{{\large\bfseries\boldmath #1}}
\newcommand{\ReasonStyle}[1]{{\normalsize #1}}

% ===== Show / Hide Reasoning Toggle =====
\newif\ifshowreason
\showreasontrue      % Instructor version
% \showreasonfalse   % Student version

% ===== Choice Item Macros (Answer + Reasoning) =====
% Usage:
%   \Correct{Answer}{Reasoning}
%   \Wrong{Answer}{Reasoning}

\newcommand{\Correct}[2]{%
  \item \AnswerStyle{\textcolor{AnsRed}{#1}}%
  \ifshowreason\\[-0.2em]
    \ReasonStyle{\textcolor{AnsRed}{#2}}%
  \fi
}

\newcommand{\Wrong}[2]{%
  \item \AnswerStyle{\textcolor{WrongBlue}{#1}}%
  \ifshowreason\\[-0.2em]
    \ReasonStyle{\textcolor{WrongBlue}{#2}}%
  \fi
}

% ===== Multiple-Choice Question Block Macro =====
% Usage: \QBlock{Q1}{stem...}{<items>}{answer key}

\newcommand{\QBlock}[4]{%
  \Needspace{8\baselineskip}%
  \noindent
  \begin{minipage}{\linewidth}
    \QuestionStyle{#1.~#2}

    \vspace{0.25em}
    \begin{enumerate}
      #3
    \end{enumerate}

    \vspace{0.1em}
    \noindent\textcolor{AnsRed}{\textit{\textbf{Answer:}~#4}}

    \vspace{0.25em}
    \hrule height 0.4pt
  \end{minipage}
  \vspace{0.75em}
}

% ===== Open-Ended Question Block Macro =====
% Usage: \OEBlock{Q1}{stem...}{answer text}

\newcommand{\OEBlock}[3]{%
  \Needspace{5\baselineskip}%
  \noindent
  \begin{minipage}{\linewidth}
    \QuestionStyle{#1.~#2}

    \vspace{0.25em}
    \noindent\textcolor{AnsRed}{\textit{\textbf{Answer:}~#3}}

    \vspace{0.25em}
    \hrule height 0.4pt
  \end{minipage}
  \vspace{0.75em}
}

% ===== Multi-Part Question Block Macro =====
% Usage: \MPBlock{Q1}{stem...}{<parts>}

\newcommand{\MPBlock}[3]{%
  \Needspace{8\baselineskip}%
  \noindent
  \begin{minipage}{\linewidth}
    \QuestionStyle{#1.~#2}

    \vspace{0.25em}
    #3

    \vspace{0.25em}
    \hrule height 0.4pt
  \end{minipage}
  \vspace{0.75em}
}

% ===== MC Sub-Part Macro =====
% Usage: \MCPart{a}{stem...}{<items>}{answer key}

\newcommand{\MCPart}[4]{%
  \par\vspace{0.35em}
  \noindent\textbf{(#1)}~#2
  \begin{enumerate}
    #3
  \end{enumerate}
  \vspace{0.1em}
  \noindent\textit{\textbf{Answer:}~#4}
  \par
}

% ===== Open-Ended Sub-Part Macro =====
% Usage: \OEPart{a}{stem...}{answer text}

\newcommand{\OEPart}[3]{%
  \par\vspace{0.35em}
  \noindent\textbf{(#1)}~#2
  \par\vspace{0.15em}
  \noindent\textit{\textbf{Answer:}~#3}
  \par
}

% ===== Skipped Question Note Macro =====
% Usage: \SkipBlock{Q1}{reason}

\newcommand{\SkipBlock}[2]{%
  \noindent\textcolor{SkipGray}{\textit{[#1 omitted: #2]}}%
  \vspace{0.5em}\par
}

% ===== Document Start =====

\begin{document}

\begin{center}
{\huge \textbf{Electromagnetism Isomorphic Questions}}\\[0.25em]
{\Large Physics Question Set}\\[0.25em]
{\large Questions generated below. Multiple-choice answers in RED; incorrect answers in BLUE.}
\end{center}

\vspace{0.75em}

\QBlock{Q1}{Two point charges are placed in vacuum. A negative charge $-q_1$ is fixed at the origin, and a positive charge $+q_2$ is located on the negative $y$-axis. What is the direction of the Coulomb force acting on $+q_2$, and what is the direction of the unit vector used in the standard form of Coulomb's law for that force?}{
\Wrong{Force on $+q_2$ points in the $-y$ direction; unit vector points in the $-y$ direction}{This treats unlike charges as if they repel, but opposite charges attract.}
\Wrong{Force on $+q_2$ points in the $+y$ direction; unit vector points in the $+y$ direction}{This matches the attraction direction for the force but incorrectly makes the unit vector point toward the source instead of from the source charge to the field point.}
\Wrong{Force on $+q_2$ points in the $-y$ direction; unit vector points in the $+y$ direction}{This reverses both the force direction and the usual source-to-field-point direction of $\hat r$.}
\Correct{Force on $+q_2$ points in the $+y$ direction; unit vector points in the $-y$ direction}{Opposite charges attract, so the force on $+q_2$ points toward the origin, which is in the $+y$ direction. In the form $\vec F = k\,q_1q_2\,\hat r/r^2$, the unit vector $\hat r$ from the source charge at the origin to the positive charge on the negative $y$-axis points in the $-y$ direction; because $q_1q_2<0$, the force is opposite to $\hat r$.}
}{D}

\QBlock{Q2}{A negative point charge $-Q$ is fixed at the origin. A positive point charge $+Q$ is placed on the negative $x$-axis. Which statement correctly describes the direction of the electric field at that point and the direction of the electric force on this positive charge?}{
\Correct{Electric field in the $+x$ direction; electric force in the $+x$ direction}{The electric field due to a negative source charge points toward the source. At a point on the negative $x$-axis, that means the field points toward the origin, that is, in the $+x$ direction; since $\vec F=q\vec E$ and the placed charge is positive, the force is in the same direction as the field.}
\Wrong{Electric field in the $+x$ direction; electric force in the $-x$ direction}{This incorrectly reverses the force direction as if the placed charge were negative rather than positive.}
\Wrong{Electric field in the $-x$ direction; electric force in the $+x$ direction}{This uses the field direction for a positive source charge instead of a negative source charge.}
\Wrong{Electric field in the $-x$ direction; electric force in the $-x$ direction}{This gets the field direction wrong and then carries that mistake into the force direction.}
}{A}

\QBlock{Q3}{A point charge is located at the center of a regular tetrahedron that acts as a Gaussian surface. Which statement is correct?}{
\Wrong{Electric flux is zero through each face}{A nonzero enclosed charge gives a nonzero total flux by Gauss's law, so the flux through each face cannot all be zero.}
\Wrong{The magnitude of the electric field is the same everywhere on each face}{Although the faces are symmetry-equivalent as whole surfaces, points on a given face are at different distances from the central charge, so the field magnitude is not uniform over the face.}
\Correct{Electric flux is the same through each face}{Gauss's law gives the total flux as $\Phi_E=q_{\text{enc}}/\varepsilon_0$. Because the charge is at the center of a regular tetrahedron, all four faces are related by symmetry, so the total flux is shared equally and each face has flux $\Phi_E/4$.}
\Wrong{None of the above is correct}{This ignores the symmetry argument that makes the flux equal through all faces.}
}{C}

\QBlock{Q4}{A very long straight wire carries a uniform positive linear charge density. A closed cylindrical Gaussian surface is placed coaxially around the wire. Which part or parts of the surface have nonzero electric flux?}{
\Wrong{Only the two flat circular end caps}{This treats the electric field as if it were along the axis, but the field of a line charge is radial.}
\Wrong{Both the curved surface and the end caps}{The radial electric field is perpendicular to the end-cap normals, so the flux through the end caps is zero.}
\Correct{Only the curved side surface}{For a long line charge, the field is radial with magnitude $E=\lambda/(2\pi\varepsilon_0 r)$. The outward area vector on the curved side is also radial, so $\vec E\cdot d\vec A\neq 0$ there, while the end-cap area vectors are along the cylinder axis and give $\vec E\cdot d\vec A=0$.}
\Wrong{No part of the surface}{If no part had flux, the net flux would be zero, contradicting Gauss's law for a surface enclosing nonzero charge.}
}{C}

\QBlock{Q5}{A positive point charge $+2q$ and a negative point charge $-q$ are separated by a distance $d$ in vacuum. The distance between them is then increased to $3d$. How does the electric potential energy of the two-charge system change?}{
\Wrong{It decreases, becoming more negative}{This ignores that for opposite charges $U=kq_1q_2/r$ is negative and becomes less negative when the separation increases.}
\Correct{It increases, becoming less negative}{The electrostatic potential energy is $U=kq_1q_2/r$. Since $q_1q_2=-2q^2<0$, increasing the separation from $d$ to $3d$ changes the energy from $U_i=-2kq^2/d$ to $U_f=-2kq^2/(3d)$; the final value is closer to zero, so the potential energy increases.}
\Wrong{It becomes exactly zero}{The energy approaches zero only at infinite separation, not merely at three times the original distance.}
\Wrong{It remains unchanged because the charges are unchanged}{Potential energy depends on both the charges and their separation, so changing $r$ changes $U$.}
}{B}

\QBlock{Q6}{An ion travels from point P to point Q in a uniform electric field. Which quantity determines the electric potential difference between P and Q?}{
\Wrong{The total length of the path between P and Q}{Electric potential difference in an electrostatic field is not determined by path length because the field is conservative.}
\Wrong{The kinetic energy of the ion}{The particle's energy may change because of the potential difference, but it does not determine the potential difference itself.}
\Wrong{The component of the displacement from P to Q perpendicular to the field}{Only the component parallel to the field enters the dot product $\Delta V=-\vec E\cdot\Delta\vec r$; perpendicular displacement does not change the potential.}
\Correct{The component of the displacement from P to Q along the electric field}{In a uniform electric field, the potential difference is $\Delta V=-\vec E\cdot\Delta\vec r$. Therefore the change in potential depends only on the component of displacement parallel or antiparallel to $\vec E$, not on the particular path taken or on the particle's properties.}
}{D}

\QBlock{Q7}{A parallel-plate capacitor has plate area $A$ and plate separation $d$. If the plate separation is tripled while the plate area remains unchanged, the capacitance:}{
\Wrong{Triples}{This reverses the relationship; capacitance is inversely proportional to plate separation.}
\Correct{Becomes one-third of its original value}{For a parallel-plate capacitor, $C=\varepsilon A/d$. Replacing $d$ with $3d$ gives $C' = \varepsilon A/(3d)=C/3$, so the capacitance is reduced to one-third of its initial value.}
\Wrong{Stays the same}{This ignores the explicit dependence of $C$ on the plate spacing.}
\Wrong{Becomes one-ninth of its original value}{This incorrectly assumes a square dependence on separation instead of the correct linear inverse dependence.}
}{B}

\QBlock{Q8}{Two point charges are separated by $16.0\,\text{cm}$ and have charges $+3.50\,\mu\text{C}$ and $-3.50\,\mu\text{C}$, respectively. What is the magnitude of the electric field at the point midway between the two charges? $k=8.99\times10^9\,\text{N}\cdot\text{m}^2/\text{C}^2$}{
\Wrong{$1.97\times10^7\,\text{N/C}$}{This likely doubles the already-combined field one extra time or uses an incorrect midpoint distance.}
\Wrong{$4.92\times10^6\,\text{N/C}$}{This is the field due to only one charge and neglects that the two contributions add at the midpoint.}
\Correct{$9.83\times10^6\,\text{N/C}$}{The midpoint is $0.080\,\text{m}$ from each charge, so the field from one charge is $E_1=k|q|/r^2=(8.99\times10^9)(3.50\times10^{-6})/(0.080)^2\approx4.92\times10^6\,\text{N/C}$. At the midpoint, the field from the positive charge points away from it and the field from the negative charge points toward it, so the two fields are in the same direction and add to give $E\approx 2E_1=9.83\times10^6\,\text{N/C}$.}
\Wrong{$0$}{This would be true for equal charges of the same sign at the midpoint, but for opposite charges the midpoint fields reinforce rather than cancel.}
}{C}

\QBlock{Q9}{A total charge of $+8.00\,\mu\text{C}$ is uniformly distributed throughout an insulating sphere of radius $0.200\,\text{m}$. What is the magnitude and direction of the electric field at a point $0.100\,\text{m}$ outside the surface of the sphere?}{
\Correct{$7.99\times10^5\,\text{N/C}$, radially outward}{Because the point is outside a spherically symmetric charge distribution, the sphere acts like a point charge located at its center. The distance from the center is $r=0.200+0.100=0.300\,\text{m}$, so $E=kQ/r^2=(8.99\times10^9)(8.00\times10^{-6})/(0.300)^2\approx7.99\times10^5\,\text{N/C}$; since the total charge is positive, the field points radially outward.}
\Wrong{$7.99\times10^5\,\text{N/C}$, radially inward}{The magnitude is correct, but the direction is wrong because the field of a positive charge points outward, not inward.}
\Wrong{$1.80\times10^6\,\text{N/C}$, radially outward}{This likely uses the wrong distance from the center, such as the distance outside the surface instead of the full radius to the point.}
\Wrong{$3.60\times10^5\,\text{N/C}$, radially outward}{This underestimates the field, typically from an algebra or distance-substitution error in $E=kQ/r^2$.}
}{A}

\QBlock{Q10}{A particle with charge $-Q$ is placed at the center of an initially uncharged conducting hollow sphere. What charges are induced on the inner and outer surfaces of the sphere?}{
\Wrong{Inner: $-Q$; outer: $+Q$}{This reverses the induced charges; the inner surface must have the opposite sign to the enclosed charge.}
\Wrong{Inner: $0$; outer: $-Q$}{This ignores the need for induced charge on the inner surface to make the electric field zero inside the conducting material.}
\Correct{Inner: $+Q$; outer: $-Q$}{To keep the electric field zero inside the conductor, the inner surface must acquire charge $+Q$, opposite to the enclosed charge $-Q$. Because the conductor as a whole was initially neutral, charge conservation requires the outer surface to carry $-Q$.}
\Wrong{Inner: $+Q$; outer: $0$}{This gets the inner surface right but violates overall charge conservation for a conductor that started neutral.}
}{C}

\QBlock{Q11}{The electric potential along the $x$-axis is given by $V(x)=4x^3-6x$, where $V$ is in volts and $x$ is in meters. What is the electric field $\vec E$ at that point?}{
\Wrong{$\vec E=(12x^2-6)\,\hat i$}{This is $dV/dx$ and misses the minus sign in the relation $\vec E=-\nabla V$.}
\Correct{$\vec E=(-12x^2+6)\,\hat i$}{In one dimension, the field is related to the potential by $E_x=-\,dV/dx$. Differentiating $V(x)=4x^3-6x$ gives $dV/dx=12x^2-6$, so $\vec E=(-12x^2+6)\hat i$.}
\Wrong{$\vec E=(x^4-3x^2)\,\hat i$}{This changes the powers as if the potential were integrated or manipulated algebraically instead of differentiated.}
\Wrong{$\vec E=(-x^4+3x^2)\,\hat i$}{This combines the wrong power-law operation with the minus sign, rather than applying $E_x=-dV/dx$ correctly.}
}{B}

\QBlock{Q12}{A $12\,\text{V}$ battery is connected in series with a $5\,\mu\text{F}$ capacitor, a $4\,\Omega$ resistor, an ammeter, and a switch that is initially open. A voltmeter is connected across the capacitor. At the instant the switch is closed, what are the ammeter and voltmeter readings, respectively?}{
\Wrong{$0\,\text{A},\ 12\,\text{V}$}{This would describe a capacitor that is already fully charged, not one that starts uncharged at the instant of switching.}
\Wrong{$0\,\text{A},\ 0\,\text{V}$}{The initial capacitor voltage is correct, but the current is not zero because an uncharged capacitor initially behaves like a short circuit.}
\Wrong{$3\,\text{A},\ 12\,\text{V}$}{This incorrectly assigns the full battery voltage across the capacitor at the same instant that it still behaves like an uncharged capacitor.}
\Correct{$3\,\text{A},\ 0\,\text{V}$}{At $t=0$, the capacitor is initially uncharged, so its voltage is $V_C(0)=0$. The entire battery voltage is therefore across the resistor, giving the initial current $I(0)=V/R=12/4=3\,\text{A}$, so the ammeter reads $3\,\text{A}$ and the voltmeter across the capacitor reads $0\,\text{V}$.}
}{D}

\QBlock{Q13}{A $9.0\,\text{V}$ battery with negligible internal resistance is used to charge a $2.5\,\text{mF}$ capacitor through a $240\,\Omega$ resistor. What is the time required for the capacitor charge to reach $70\%$ of its maximum value?}{
\Wrong{$0.18\,\text{s}$}{This is too small and usually comes from using the wrong fraction inside the logarithm or omitting the factor $RC$.}
\Wrong{$0.42\,\text{s}$}{This can result from using $-RC\ln(0.50)$ or another incorrect target fraction rather than $0.70$.}
\Correct{$0.72\,\text{s}$}{For charging, $q(t)=Q_{\max}(1-e^{-t/RC})$. Setting $q/Q_{\max}=0.70$ gives $e^{-t/RC}=0.30$, so $t=-RC\ln(0.30)$; with $RC=(240)(2.5\times10^{-3})=0.600\,\text{s}$, the time is $t\approx0.600\ln(1/0.30)\approx0.72\,\text{s}$.}
\Wrong{$1.20\,\text{s}$}{This overestimates the time, as if the capacitor needed about two time constants to reach $70\%$, which is not consistent with the exponential charging law.}
}{C}

\QBlock{Q14}{At one instant a proton $(q=+1.60\times10^{-19}\,\text{C})$ is moving in the $xz$-plane with velocity components $v_x=2.0\times10^5\,\text{m/s}$ and $v_z=4.0\times10^5\,\text{m/s}$. A magnetic field of $0.50\,\text{T}$ is in the positive $x$-direction. At that instant, the magnetic force on the proton is:}{
\Correct{$(3.2\times10^{-14}\,\text{N})\,\hat j$}{The magnetic force is $\vec F=q\,\vec v\times\vec B$. Only the $v_z$ component contributes because $v_x$ is parallel to $\vec B$; thus $|\vec F|=qv_zB=(1.60\times10^{-19})(4.0\times10^5)(0.50)=3.2\times10^{-14}\,\text{N}$, and since $\hat k\times\hat i=\hat j$ for a positive charge, the direction is $+\hat j$.}
\Wrong{$(-3.2\times10^{-14}\,\text{N})\,\hat j$}{This gets the magnitude right but reverses the direction, as if the particle had negative charge instead of positive charge.}
\Wrong{$(1.6\times10^{-14}\,\text{N})\,\hat j$}{This uses only half the correct magnitude, typically from substituting the wrong velocity component or field value.}
\Wrong{$(4.5\times10^{-14}\,\text{N})\,\hat j$}{This likely combines both velocity components as if all of $\vec v$ were perpendicular to $\vec B$, which is not true here.}
}{A}

\end{document}
